\documentclass[10pt,a4paper]{article}
\usepackage{amsmath,amssymb,amsthm,graphicx}
\usepackage{hyperref}
\usepackage{latexsym}
\usepackage{amssymb}
\usepackage{amsthm}
\usepackage{ulem}
\usepackage{enumitem}
\setlist{itemsep=0.4pt,topsep=0pt}
\parindent 0pt
\parindent 0pt
\parskip 12pt
\newcommand{\HCF}{\mathrm{hcf}}
\newcommand{\ZZ}{\mathbb{Z}}
\newcommand{\rk}{\mathrm{rk}}
\DeclareMathOperator{\len}{length}
\DeclareMathOperator{\bl}{b}
\DeclareMathOperator{\cl}{cl}
\DeclareMathOperator{\di}{d}
\DeclareMathOperator{\Cl}{Cl}
\DeclareMathOperator{\acl}{acl}
\DeclareMathOperator{\tp}{tp}
\DeclareMathOperator{\Th}{Th}
\newcommand{\Cor}{\mathcal{C}_\mathrm{or}}
\newcommand{\Cun}{\mathcal{C}_\mathrm{un}}
\newcommand{\mons}{\mathbb{M}}
\newcommand{\strong}{\mathrm{strong}}
\newcommand{\axioms}{\Delta_{\mathrm{FSTS}}}
\newtheoremstyle{note}% hnamei
{8pt}% hSpace abovei
{14pt}% hSpace belowi
{\normalfont}% hBody fonti
{}% hIndent amounti1
{\bfseries}% hTheorem head fonti
{ }% hPunctuation after theorem headi
{.5em}% hSpace after theorem headi2
{}% hTheorem head spec (can be left empty, meaning ‘normal’)i
\theoremstyle{note}
\newtheorem{thm}{Theorem}[section]
%\newtheorem{theorem}{Theorem}
\newtheorem{lemma}[thm]{Lemma} 
\newtheorem{df}[thm]{Definition}
\newtheorem{rmk}[thm]{Remark} 
\newtheorem{prop}[thm]{Proposition}
\newtheorem{fact}[thm]{Fact}
\newtheorem{cor}[thm]{Corollary}
\newtheorem{conj}[thm]{Conjecture}
\newtheorem{ex}[thm]{Example}
\newtheorem{qu}[thm]{Question}
\newtheorem*{notation}{Notation} 
\title{The diamond conditions for unconfined systems and a minimal pair}
\author{Silvia Barbina, Enrique Casanovas}
\date{arxiv:2411.13723v3; Creative Commons Attribution 4.0 International}
\begin{document}
\maketitle
\noindent\textit{Editorial note.} The counted statement is the author's environment at source0.tex lines 650-655 and its proof environment opens at line 657; the pairing is by adjacency and the retained proof range is the whole environment.\par\medskip
\noindent\textit{Editorial note.} The preamble is the author's own package lines, macro block and theorem-environment declarations (source lines 4-65); the author's optioned amsmath line and the fleqn-dependent layout line are not retained.\par\medskip
\noindent\textit{Editorial note.} The source carries no label-number dump and no counter values are imposed, so the printed numbering is the wrapper's default; identity is the author label, or the pinned source position where the author gave none.\par\medskip
\noindent\textit{Editorial note.} Metadata provenance: the source record of the pinned copy and the licence evidence file are both present; the licence is CC BY 4.0 with status supported, and the card's source\_version matches the value recorded in that evidence file.\par\medskip
\noindent\textit{Editorial note.} r002 note: r001 failed because the retained author proof cites two bibliography keys of the article; r002 bundles the article's own bibliography file and changes nothing else.\par\medskip
\noindent\textit{Editorial scope.} One original result of the pinned author manuscript, counted once: the two diamond conditions for unconfined systems (author label diamond).\par\medskip
\noindent\textit{Allowed standard prerequisites (editorial).}\begin{itemize}
\item HF-ordering. A linear ordering in which every element has at most one pair of smaller elements forming a block with it.. It is what the cited proposition supplies for an unconfined system.
\item Elementary substructure. A\$\textbackslash{}leq\$B means that A is an elementary substructure of B, so that the closure conditions transfer.. It is the relation appearing in the second item.
\item Inherited ordering. A subset of an ordered structure inherits the ordering, and the initial segment condition restricts to it.. It is the mechanism of the second item.
\end{itemize}\medskip
	\begin{prop}\label{unconfined} \begin{enumerate}
			\item Free Steiner triple systems are unconfined.
%			\item If $A \subseteq M$ is a finite subset of the Steiner quasigroup $M$, then $A$ is contained in a finitely generated free quasigroup. {\color{red} move  later}
			\item Finitely generated unconfined STSs are free.
			\item A partial STS is unconfined if and only if it has an HF-ordering.
		\end{enumerate}
	\end{prop}

\begin{lemma}\label{diamond} Let $A,A_1,A_2 \in \Cun$. Then
	\begin{enumerate}
		\item $\emptyset \leq A$  
	\item If $B\in \Cun$ is such that $A_1\subseteq B$ and $A_2\leq B$, then $A_1\cap A_2\leq A_1$.
	\end{enumerate}
\end{lemma}

\begin{proof} 1. By Proposition~\ref{unconfined},  since every $A\in \Cun$ has an HF-ordering.
	
	2. Given an HF-ordering of $B$ with initial segment $A_2$, the elements of $A_1$ inherit an HF-ordering with initial segment $A_1\cap A_2$.
\end{proof}

\begin{thebibliography}{9}
\bibitem{barcas2}
S.~Barbina and E.~Casanovas.
\newblock Endomorphisms of free {S}teiner quasigroups, 2024.
\newblock{\url{arXiv:2412.06055}}.

\bibitem{casanovas}
E.~Casanovas.
\newblock Compactly expandable models and stability.
\newblock {\em The Journal of Symbolic Logic}, 60:673--683, 1995.
\bibitem{ganter1}
B.~Ganter and H.~Werner.
\newblock Equational classes of {S}teiner systems.
\newblock {\em Algebra Universalis}, 5:125--140, 1975.

\bibitem{ganter2}
B.~Ganter and H.~Werner.
\newblock Co-ordinatizing {S}teiner systems.
\newblock In C.C. Linder and A.~Rosa, editors, {\em Topics on {S}teiner
  systems}, pages 3--24. North Holland, 1980.
\end{thebibliography}
\end{document}
